% answers_short.tex —— 2017 年·北京卷·理科 · 简短答案速查
% 由 tools/make_short_answers.py 自 answers.tex 的 \gksolhead/\gkanswer 提取，
% 并与 2017/2017_数学真题_北京理_答案要点.md 逐题比对后生成，勿手改（改源 MD 后重跑）。

\begin{gknotice}
\textbf{使用说明}：本节只给最终答案便于快速核对，完整推导见解析版。答案为编写者独立推导，\textbf{非官方标准答案}。
\end{gknotice}

\gkpart{选择题}
\begin{center}\begin{tabular}{ccccccccc}
\toprule
题号 & 1 & 2 & 3 & 4 & 5 & 6 & 7 & 8 \\
\midrule
答案 & A & B & C & D & A & A & B & D \\
\bottomrule
\end{tabular}\end{center}

\gkpart{填空题}
\begin{itemize}[leftmargin=2.6em]
  \item[\textbf{9.}] $2$
  \item[\textbf{10.}] $1$
  \item[\textbf{11.}] $1$
  \item[\textbf{12.}] $-7/9$
  \item[\textbf{13.}] 例如 $(-1,-2,-3)$，答案不唯一
  \item[\textbf{14.}] (1) $Q_1$；(2) $p_2$
\end{itemize}

\gkpart{解答题}
\begin{itemize}[leftmargin=2.6em]
  \item[\textbf{15.}] (1) $3\sqrt3/14$；(2) $6\sqrt3$
  \item[\textbf{16.}] (1) $M$ 为 $PB$ 中点；(2) $60^\circ$；(3) $2\sqrt6/9$
  \item[\textbf{17.}] (1) $3/10$；(2) $P(\xi=0,1,2)=(1/6,2/3,1/6)$，$E(\xi)=1$；(3) 服药者的 $y$ 方差较大
  \item[\textbf{18.}] (1) $y^2=x$，焦点 $(1/4,0)$，准线 $x=-1/4$；(2) $A$ 为 $BM$ 中点
  \item[\textbf{19.}] (1) $y=1$；(2) 最大值 $1$，最小值 $-\pi/2$
  \item[\textbf{20.}] (1) $c_1=0,c_2=-1,c_3=-2,c_n=1-n$；(2) 按 $\{a_n\}$ 公差的正、负、零分类证明
\end{itemize}
